\section{The micro-kernels}
\label{sec:mk}

This section describes how the components of Section~\ref{sec:components} are actually evaluated:
the staging, the loop structure, the compile-time dispatch, and what differs per scheme and per
layout. All counts below are properties of the written kernel and can be checked by reading it;
timings belong to the paper.

\subsection{Lane blocking and the packs}
\label{sec:mk:staging}

Every kernel processes a \emph{lane group} of elements at once. The width is
\code{CVFEM\_HEX8\_VEC\_SIZE = VEC\_BYTES / sizeof(scalar\_t)}, with \code{VEC\_BYTES} defaulting
to 128, so sixteen elements in double precision. Element data is staged into lane-major structs,
indexed \code{[node][lane]} --- or \code{[node][component][lane]} for the gradient packs --- so
that a fixed node index over varying lane is contiguous and every read in the surface loop is a
unit-stride vector load rather than a gather:

\begin{center}
\begin{tabular}{@{}lll@{}}
\toprule
pack & contents & read by \\
\midrule
\code{Hex8InputPack}    & $u_x,u_y,u_z,p$ per node            & every scheme \\
\code{Hex8ResidualPack} & the four accumulators per node      & every scheme \\
\code{Hex8RhieChowPack} & $\nabla p$ per node; $G_\scs$, $w_\scs$ per surface & Rhie--Chow arms \\
\code{Hex8UGradPack}    & $\nabla\uvec$ (nine components) and coordinates per node & higher-order arms \\
\bottomrule
\end{tabular}
\end{center}

\code{Hex8RhieChowPack} notably does \emph{not} carry node coordinates. It used to, purely to form
$\dvec$ by differencing; since the affine path takes $\dvec$ from the Jacobian
(Section~\ref{sec:geom:edges}) those three arrays have no reader, and removing them took 3\,KB of
stack per lane group and twenty-four indexed gathers per element out of the sweep. The
isoparametric kernels keep their own \code{Hex8CoordPack}.

\subsection{Loop structure}
\label{sec:mk:structure}

The residual's twelve surfaces are evaluated in \emph{one} lane loop
(\code{cvfem\_hex8\_conv\_all\_simd}), with each surface's contribution accumulated into
thirty-two lane-private scalars --- eight nodes times four rows --- and each stored once at the
end. The alternative, a vectorised lane loop per surface, writes a node's four accumulators three
times over, once per direction group, and measured slower.

The binding constraint on this loop is register capacity, and it is worth stating plainly because
it bounds what any rescheduling can achieve. One vector iteration holds thirty-two accumulators,
and each surface needs roughly two dozen more live values --- the area vector, both nodes'
velocity and pressure, and for the Rhie--Chow arms the pressure gradients and the two tabulated
surface scalars. That is on the order of fifty-six live vector values against the thirty-two
registers AArch64 provides, so the register allocator spills regardless of the order the surfaces
are visited in. Measured on the shipping binary, roughly $37\%$ of the convective sweep's
instructions were spill traffic against a $1952$-byte frame.

Three reschedulings were built and measured against this shape, and the pattern is instructive:
varying the lane width did nothing (the live set is per vector iteration, not across unrolled
iterations); fusing the surfaces as described above gained a few percent; and splitting the
momentum and continuity rows into separate passes --- which halves the accumulator count, exactly
the lever the spill count suggests --- \emph{lost} a third, because the per-surface inputs already
exceed what the accumulators leave free and the split only duplicated the mass-flux work. The
conclusion carried into Section~\ref{sec:geom:edges} is that the live set, not the schedule, is the
constraint, and that shrinking it means changing what a surface \emph{reads}.

\subsection{Compile-time dispatch}

Everything that is uniform over a sweep is a template parameter, not a runtime argument:

\begin{center}
\begin{tabular}{@{}ll@{}}
\toprule
flag & selects \\
\midrule
\code{RC}  & the Rhie--Chow term \\
\code{QG}  & the direction's pressure-gradient term in the Jacobian \\
\code{EPS} & the Harten band \\
\code{HO}  & the deferred correction \\
\code{LIM} & which limiter, $0..3$ \\
\code{NC}  & the field count in the reconstruction \\
\bottomrule
\end{tabular}
\end{center}

The reason is uniform across the file and was established by measurement rather than taste: a
sweep-uniform test left inside the vector body costs between $1.8\times$ and $2\times$, because it
tips the compiler's cost model out of vectorising the loop that contains it. A guard on the
Rhie--Chow coefficient cost $1.83\times$ this way; a runtime limiter select cost $2\times$. The
flags are dispatched once, outside the lane loop, by a macro that keeps the long argument list in
one place --- sixteen hand-written copies of it is where a transposed pack would hide.

\subsection{Per-scheme kernels}
\label{sec:mk:perscheme}

First-order upwinding runs the hand-written sweep with \code{HO=false}. The four higher-order
schemes are evaluated by \emph{generated} whole-element kernels, synthesised by
\code{python/synthesize\_cvfem\_hex8\_ns\_upwind\_sympy.py} and emitted into
\code{cvfem\_hex8\_ns\_upwind\_sympy\_kernels.hpp}: one per limiter, in two variants each with and
without Rhie--Chow, and all with a single \code{\#pragma omp simd} lane loop wrapping the entire
body:

\begin{center}
\begin{tabular}{@{}llrr@{}}
\toprule
scheme & kernel & lines, plain & lines, $+$RC \\
\midrule
unlimited          & \code{...\_defcor\_lim0\_simd} &  704 &  880 \\
bounded-face clip  & \code{...\_defcor\_lim1\_simd} &  798 &  963 \\
Venkatakrishnan    & \code{...\_defcor\_lim2\_simd} & 2290 & 2427 \\
Darwish--Moukalled & \code{...\_defcor\_lim3\_simd} & 1487 & 1611 \\
\bottomrule
\end{tabular}
\end{center}

The size ordering follows the limiter's algebra: the clip is a pair of comparisons, van Leer's
form has one division, and Venkatakrishnan's rational function with its deactivation term is three
times the body of either. Common subexpression elimination is applied across the whole element, so
the twelve surfaces share the centroid and geometry work that twelve separately compiled
per-surface functions could not.

These kernels contain no \code{if} statements at all --- the limiter logic is entirely ternaries,
which lower to selects and vectorise. That property is load-bearing and was not free: a ternary
whose \emph{condition} is itself a ternary does \emph{not} vectorise, because the compiler
if-converts the inner one into a PHI node and then cannot select over it. Venkatakrishnan's guard
had that shape (Section~\ref{sec:schemes}), and the kernel issued zero vector instructions until
the guard was moved onto the increment. The lesson generalises past this one kernel: in generated
code the property to check is not "is there a lane loop" --- all four always had one --- but
whether the compiler acted on it, and only a disassembly or a vectoriser report answers that.

\subsection{What the layout changes}

The two layouts differ in how an element's contribution reaches the global vector, not in the
arithmetic. The \emph{packed} layout gathers pack-locally, accumulates into pack-local
indices, and closes through a two-phase deterministic ghost reduction; the \emph{standard}
layout scatters with atomics. Consequences worth knowing when reading measurements:

\begin{itemize}
\item The packed layout is bitwise reproducible by construction. The standard layout is not ---
      its fingerprint varies run to run with the same binary, because the atomic order varies ---
      so a fingerprint comparison is only meaningful within the packed layout.
\item The packed layout spends a larger share of its time in the surface loop, so changes to the
      per-surface work show up there more strongly, and changes to per-element work (staging,
      gathers) show up more strongly on the standard layout.
\item Checksums are near-cancelling sums of order $10^{-14}$ built from terms of order one, so a
      last-bit difference reads as a large relative change. They are not a usable discriminator;
      the per-node cross-layout comparisons are.
\end{itemize}
